% ========== GENERAL MATH COMMANDS ==========
% vocab:
\newcommand{\vocab}[1]{\textbf{\color{black!90}\boldmath #1}}

% inverse trigs: \arcsec, \arccot, \arccsc
\DeclareMathOperator{\arcsec}{arcsec}
\DeclareMathOperator{\arccot}{arccot}
\DeclareMathOperator{\arccsc}{arccsc}

% contradiction
\newcommand{\contradiction}{
  {\hbox{
    \setbox0=\hbox{\(\mkern-3mu{\times}\mkern-3mu\)}
    \setbox1=\hbox to0pt{\hss\copy0\hss}
    \copy0\raisebox{0.5\wd0}{\copy1}\raisebox{-0.5\wd0}{\box1}\box0}
  }
}

% better looking mod:
\renewcommand{\mod}[1]{\ \text{mod}\ #1}

% \abs and \norm
\DeclarePairedDelimiter\abs{\lvert}{\rvert}
\DeclarePairedDelimiter\norm{\lVert}{\rVert}
% floor and ceil
\DeclarePairedDelimiter\ceil{\lceil}{\rceil}
\DeclarePairedDelimiter\floor{\lfloor}{\rfloor}

% === shortcuts for mathbb, mathcal, and mathscr ===
\newcommand{\NN}{\mathds{N}}
\newcommand{\ZZ}{\mathds{Z}}
\newcommand{\QQ}{\mathds{Q}}
\newcommand{\RR}{\mathds{R}}
\newcommand{\CC}{\mathds{C}}
\newcommand{\PP}{\mathds{P}}
\newcommand{\FF}{\mathds{F}}

\ExplSyntaxOn
% \bX -> \mathds{X}
\clist_map_inline:nn {N,Z,Q,R,C,P,F,S,T}
{
  \cs_new:cpn { b#1 } { \mathds{#1} }
}
% \cX -> \mathcal{X}
\clist_map_inline:nn {A,B,C,D,E,F,G,H,I,J,K,L,M,N,O,P,Q,R,S,T,U,V,W,X,Y,Z}
{
  \cs_new:cpn { c#1 } { \mathcal{#1} }
}
% \sX -> \mathscr{X}
\clist_map_inline:nn {A,B,C,D,E,F,K,L,M,P,R,S,T,U,V}
{
  \cs_new:cpn { s#1 } { \mathscr{#1} }
}
\ExplSyntaxOff

\newcommand{\fM}{\mathfrak{M}}

% ==================== MATH, BY AREA ====================
\DeclareMathOperator*{\argmax}{arg\,max}
\DeclareMathOperator*{\argmin}{arg\,min}
% ========== Set Theory ==========
% complement: \complement
\renewcommand{\complement}[0]{\mathsf{c}}
% nicer emptyset
\let\emptyset\varnothing
% cardinality: \card
\newcommand{\card}[1]{\left|#1\right|}
\newcommand{\dom}{\mathcal{D}}
% indicator function
\makeatletter
\NewDocumentCommand{\ind}{}{%
  \mathds{1}%
  \@ifnextchar\bgroup{\ind@witharg}{}%
}
\NewDocumentCommand{\ind@witharg}{m}{%
  \!\left\{\,#1\,\right\}%
}
\makeatother

% ========== Point Set Topology ==========
\DeclareMathOperator{\Int}{Int}

% ========== Probability ==========
\DeclareMathOperator{\Var}{\mathds{V}}
\DeclareMathOperator{\sd}{sd}
\DeclareMathOperator{\se}{se}
\DeclareMathOperator*{\plim}{plim}
\newcommand{\Cov}{\mathds{C}}
\DeclareMathOperator{\Corr}{Corr}
\renewcommand{\Pr}{\mathds{P}}
\renewcommand{\P}{\mathds{P}}
\newcommand{\E}{\mathds{E}}
\DeclareMathOperator{\MSE}{MSE}
\DeclareMathOperator{\Bias}{Bias}

\newcommand{\disteq}{\stackrel{\mathrm{d}}{=}}
\newcommand{\distto}{\xrightarrow{\mathrm{d}}}
\newcommand{\longdistto}{\xrightarrow{\ \mathrm{d}\ }}
\newcommand{\probto}{\xrightarrow{\mathrm{p}}}
\newcommand{\longprobto}{\xrightarrow{\ \mathrm{p}\ }}

\newcommand{\iidsim}{\overset{\mathrm{iid}}{\sim}}

\DeclareMathOperator{\Bernoulli}{Bernoulli}
\DeclareMathOperator{\Binomial}{Binomial}
\DeclareMathOperator{\Multinomial}{Multinomial}
\DeclareMathOperator{\Geometric}{Geometric}
\DeclareMathOperator{\Poisson}{Poisson}
\DeclareMathOperator{\Exponential}{Exponential}
\DeclareMathOperator{\DirichletDist}{Dirichlet}
\DeclareMathOperator{\GammaDist}{Gamma}
\DeclareMathOperator{\BetaDist}{Beta}
\DeclareMathOperator{\Normal}{\mathcal{N}}
\DeclareMathOperator{\Uniform}{Uniform}
\let\Unif\Uniform

% ========== Linear Algebra ==========
\DeclareMathOperator{\Id}{Id}
\DeclareMathOperator{\tr}{tr}
\DeclareMathOperator{\rank}{rank}
\DeclareMathOperator{\RREF}{RREF}
\DeclareMathOperator{\sgn}{sgn}
\DeclareMathOperator{\Span}{span}

\newcommand\T{{\mathpalette\raiseT\intercal}}
\newcommand\raiseT[2]{\raisebox{0.25ex}{$#1#2$}}
\newcommand*{\tran}{\T} % deprecated



\newcommand\spanset[1]{\ensuremath\Span\left(#1\right)}

% nicer vector
\renewcommand{\vec}[1]{\boldsymbol{\mathbf{#1}}}

% matrices
\newcommand{\bmat}[1]{\begin{bmatrix}#1\end{bmatrix}}
\newcommand{\vmat}[1]{\begin{vmatrix}#1\end{vmatrix}}
\newcommand{\pmat}[1]{\begin{pmatrix}#1\end{pmatrix}}

% ========== Analysis ==========
\renewcommand{\Re}{\operatorname{Re}}
\renewcommand{\Im}{\operatorname{Im}}
\newcommand*{\ii}{\mathrm{i}}
\let\im\ii % deprecated
\let\I\ii % deprecated

\newcommand{\dd}{\mathop{}\!\mathrm{d}}
\newcommand{\DD}{\mathrm{D}}

\newcommand{\ran}{\mathcal{R}}
\renewcommand{\ker}{\mathcal{N}}

\DeclareMathOperator{\dist}{dist}
\DeclareMathOperator{\diam}{diam}

\DeclareMathOperator{\graph}{G}
\let\gr\graph
\DeclareMathOperator{\epi}{epi}

\DeclareMathOperator{\supp}{supp}

\DeclareMathOperator{\Lip}{Lip}

% convex hull
\DeclareMathOperator{\conv}{conv}

% ========== Spectral Theory ==========
\DeclareMathOperator{\res}{res}
\DeclareMathOperator{\spec}{spec}

% ========== Measure Theory ==========
\newcommand{\essran}{\mathrm{ess\ ran}\ }
\newcommand{\weakto}{\rightharpoonup}

% ========== Algebra ==========
\DeclareMathOperator{\GL}{GL}

% ========== ECONOMICS ==========
% ========== Econometrics ==========
\DeclareMathOperator{\SST}{SST}
\DeclareMathOperator{\SSE}{SSE}
\DeclareMathOperator{\SSR}{SSR}
\DeclareMathOperator{\AR}{AR}
\DeclareMathOperator{\MA}{MA}
\newcommand{\indep}{\perp \!\!\! \perp}

